% R10_Isomorphism_Theorem.tex
% monogate Cost Theory — SuperBEST Isomorphism Theorem
% Proves that equations from different scientific domains share identical
% minimal SuperBEST DAG topologies, differing only in terminal labels.
% Date: 2026-04-20
% Status: THEOREM

\documentclass[12pt,a4paper]{article}
\usepackage{amsmath,amssymb,amsthm}
\usepackage{geometry}
\usepackage{booktabs}
\usepackage{array}
\usepackage{longtable}
\usepackage{xcolor}
\usepackage{hyperref}
\geometry{margin=2.5cm}

% ── Theorem environments ────────────────────────────────────────────────────
\newtheorem{theorem}{Theorem}[section]
\newtheorem{lemma}[theorem]{Lemma}
\newtheorem{corollary}[theorem]{Corollary}
\newtheorem{proposition}[theorem]{Proposition}
\theoremstyle{definition}
\newtheorem{definition}[theorem]{Definition}
\newtheorem{conjecture}[theorem]{Conjecture}
\newtheorem{remark}[theorem]{Remark}
\newtheorem{example}[theorem]{Example}
\newtheorem{observation}[theorem]{Observation}

% ── Operator macros ─────────────────────────────────────────────────────────
\newcommand{\EML}{\operatorname{EML}}
\newcommand{\EDL}{\operatorname{EDL}}
\newcommand{\EXL}{\operatorname{EXL}}
\newcommand{\EAL}{\operatorname{EAL}}
\newcommand{\EMN}{\operatorname{EMN}}
\newcommand{\DEML}{\operatorname{DEML}}
\newcommand{\ELAd}{\operatorname{ELAd}}
\newcommand{\ELSb}{\operatorname{ELSb}}
\newcommand{\LEAd}{\operatorname{LEAd}}
\newcommand{\LEdiv}{\operatorname{LEdiv}}
\newcommand{\LEprod}{\operatorname{LEprod}}
\newcommand{\EPL}{\operatorname{EPL}}
\newcommand{\DEAL}{\operatorname{DEAL}}
\newcommand{\DEXL}{\operatorname{DEXL}}
\newcommand{\DEDL}{\operatorname{DEDL}}
\newcommand{\DEPL}{\operatorname{DEPL}}
\newcommand{\DEMN}{\operatorname{DEMN}}

% ── Cost macros ──────────────────────────────────────────────────────────────
\newcommand{\Cost}{\operatorname{Cost}}
\newcommand{\NC}{\operatorname{NaiveCost}}
\newcommand{\SD}{\operatorname{SharingDiscount}}
\newcommand{\PB}{\operatorname{PatternBonus}}
\newcommand{\MinDAG}{\operatorname{MinDAG}}
\newcommand{\Fam}{\mathcal{F}}

% ── Tree-node shorthand ──────────────────────────────────────────────────────
\newcommand{\mexp}{\mathtt{exp}}
\newcommand{\mmul}{\mathtt{mul}}
\newcommand{\mdiv}{\mathtt{div}}
\newcommand{\madd}{\mathtt{add}}
\newcommand{\msub}{\mathtt{sub}}
\newcommand{\mln}{\mathtt{ln}}
\newcommand{\mrecip}{\mathtt{recip}}
\newcommand{\mneg}{\mathtt{neg}}

% ── Isomorphism symbol ───────────────────────────────────────────────────────
\newcommand{\siso}{\cong_{\mathrm{SB}}}

% ── Column types ─────────────────────────────────────────────────────────────
\newcolumntype{L}[1]{>{\raggedright\arraybackslash}p{#1}}
\newcolumntype{C}[1]{>{\centering\arraybackslash}p{#1}}

\title{The SuperBEST Isomorphism Theorem\\[6pt]
       \large Cross-Domain Equations with Identical Minimal Tree Topology}
\author{Arturo R.\ Almaguer\\
\small monogate research \quad \texttt{almaguer1986@gmail.com}}
\date{April 2026}

% ════════════════════════════════════════════════════════════════════════════
\begin{document}
\maketitle

% ── Abstract ────────────────────────────────────────────────────────────────
\begin{abstract}
We prove the \emph{SuperBEST Isomorphism Theorem}: equations from unrelated
scientific disciplines often possess identical minimal SuperBEST DAGs, differing
only in the names of their terminal nodes (free variables and constants).
We introduce the formal notion of SuperBEST-isomorphism, exhibit eight
cross-domain isomorphism families spanning physics, chemistry, pharmacokinetics,
economics, and machine learning, and show that these families define equivalence
classes of equations under the isomorphism relation.  The theorem implies a
\emph{Revealed Structure Corollary}: domain-specific notation conceals shared
mathematical infrastructure that becomes manifest only at the level of minimal
operator trees.
\end{abstract}

\tableofcontents
\bigskip

% ════════════════════════════════════════════════════════════════════════════
\section{SuperBEST Isomorphism: Definition}
\label{sec:def}
% ════════════════════════════════════════════════════════════════════════════

Throughout this paper, $\mathcal{F}_{16}$ denotes the 16-operator monogate
family and $\MinDAG(E)$ denotes the unique minimal SuperBEST v3 DAG for
expression $E$ (unique up to labelling of equivalent subexpressions).

% ── 1.1 Preliminary: minimal DAGs ───────────────────────────────────────────
\subsection{Minimal SuperBEST DAGs}

\begin{definition}[Minimal SuperBEST DAG]
\label{def:mindag}
The \emph{minimal SuperBEST DAG} of an expression $E$, written $\MinDAG(E)$,
is the directed acyclic graph with minimum node count over $\mathcal{F}_{16}$
that evaluates to $E$.  Internal nodes are labelled by operators from
$\mathcal{F}_{16}$; leaf (terminal) nodes are labelled by free variables or
numeric constants.  Node count equals $\Cost(E)$.
\end{definition}

\begin{remark}[Existence and near-uniqueness]
The SuperBEST v3 routing table guarantees existence of $\MinDAG(E)$ for every
expression $E$ computable over $\mathcal{F}_{16}$.  Structural uniqueness holds
for all expressions in the standard scientific corpus: no two structurally
distinct DAGs achieve the same node count for any equation appearing in our
catalogue.
\end{remark}

% ── 1.2 Isomorphism ─────────────────────────────────────────────────────────
\subsection{SuperBEST Isomorphism}

\begin{definition}[SuperBEST-isomorphic expressions]
\label{def:iso}
Two arithmetic expressions $E_1$ and $E_2$ are \emph{SuperBEST-isomorphic},
written $E_1 \siso E_2$, if there exists a graph bijection
\[
  \varphi \colon \mathrm{nodes}(\MinDAG(E_1))
             \;\longrightarrow\;
             \mathrm{nodes}(\MinDAG(E_2))
\]
satisfying:
\begin{enumerate}
  \item \textbf{Operator preservation.}
        For every internal node $u$ in $\MinDAG(E_1)$:
        \[
          \mathrm{op}(u) \;=\; \mathrm{op}(\varphi(u)).
        \]
  \item \textbf{Topology preservation.}
        $u$ is a child of $v$ in $\MinDAG(E_1)$ if and only if $\varphi(u)$
        is a child of $\varphi(v)$ in $\MinDAG(E_2)$.
  \item \textbf{Terminal-to-terminal mapping.}
        $u$ is a terminal (leaf) node in $\MinDAG(E_1)$ if and only if
        $\varphi(u)$ is a terminal node in $\MinDAG(E_2)$.
\end{enumerate}
The bijection $\varphi$ need not preserve terminal \emph{labels}
(variable names or constant values); terminals may be freely renamed across
the isomorphism.
\end{definition}

\begin{remark}[What isomorphism captures]
Condition~(1) states that corresponding nodes compute the same primitive
operation.  Condition~(2) states that the wiring pattern---which node feeds
into which---is identical in both DAGs.  Condition~(3) ensures that the
``input slots'' align: the same number of free inputs appear at the same
structural positions.  Together, the three conditions say that $E_1$ and
$E_2$ are the same computation up to renaming of inputs.
\end{remark}

% ── 1.3 Equivalence relation ─────────────────────────────────────────────────
\subsection{Isomorphism Classes}

\begin{proposition}[$\siso$ is an equivalence relation]
\label{prop:equiv}
SuperBEST-isomorphism is reflexive, symmetric, and transitive on the set of
arithmetic expressions computable over $\mathcal{F}_{16}$.
\end{proposition}

\begin{proof}
\textit{Reflexivity.} The identity map $\mathrm{id}$ on $\MinDAG(E)$ satisfies
all three conditions for $E \siso E$.

\textit{Symmetry.} If $\varphi\colon \MinDAG(E_1) \to \MinDAG(E_2)$ witnesses
$E_1 \siso E_2$, then the inverse bijection $\varphi^{-1}$ witnesses
$E_2 \siso E_1$, since all three conditions are symmetric in $E_1$ and $E_2$
when restated through $\varphi^{-1}$.

\textit{Transitivity.} If $\varphi_{12}$ witnesses $E_1 \siso E_2$ and
$\varphi_{23}$ witnesses $E_2 \siso E_3$, then
$\varphi_{23} \circ \varphi_{12}$ witnesses $E_1 \siso E_3$: composition
preserves all three conditions.
\end{proof}

\begin{definition}[Isomorphism class and canonical representative]
\label{def:class}
The \emph{SuperBEST isomorphism class} of $E$ is the set
$[E] = \{E' : E' \siso E\}$.  The \emph{canonical representative tree} of
$[E]$ is the abstract operator tree obtained from $\MinDAG(E)$ by replacing
every terminal label with a placeholder variable $x_1, x_2, \ldots$
in left-to-right, leaf-first order.
\end{definition}

% ── 1.4 Cost invariance ──────────────────────────────────────────────────────
\subsection{Cost is an Isomorphism Invariant}

\begin{proposition}[Cost invariance]
\label{prop:cost-inv}
If $E_1 \siso E_2$ then $\Cost(E_1) = \Cost(E_2)$.
\end{proposition}

\begin{proof}
The bijection $\varphi$ is a graph isomorphism between $\MinDAG(E_1)$ and
$\MinDAG(E_2)$.  In particular, $|\mathrm{nodes}(\MinDAG(E_1))| =
|\mathrm{nodes}(\MinDAG(E_2))|$.  Since $\Cost(E) = |\mathrm{nodes}(\MinDAG(E))|$
by Definition~\ref{def:mindag}, equality of cost follows immediately.
\end{proof}

% ════════════════════════════════════════════════════════════════════════════
\section{Eight Cross-Domain Isomorphism Families}
\label{sec:families}
% ════════════════════════════════════════════════════════════════════════════

We now exhibit eight concrete isomorphism families.  For each family we
(i)~give the canonical operator tree, (ii)~list the member equations from
distinct domains, (iii)~exhibit the explicit bijection between any two
representative members, and (iv)~state the cost.

% ── 2.1 Exponential growth/decay family ─────────────────────────────────────
\subsection{Family 1: Exponential Scale-and-Decay (5n)}
\label{sec:fam1}

\paragraph{Canonical tree.}
\begin{equation}
\label{eq:fam1-tree}
  T_1 \;=\; \mmul\!\bigl(a,\; \mexp(\mmul(b,\, c))\bigr)
\end{equation}
with three terminal slots $a$, $b$, $c$.  Cost computation:
$c_{\mmul} = 2$, $c_{\mexp} = 1$, $c_{\mmul} = 2$; total $= 5$.
All three terminals are live, so $\NC = 5$, $\SD = \PB = 0$, and
$\Cost(T_1) = 5$.

\paragraph{Member equations.}
\begin{align*}
  E_{1a} &: N(t) = N_0 \cdot \exp(-\lambda t)
            && \text{(Radioactive decay)} \\
  E_{1b} &: V(t) = V_0 \cdot \exp(r t)
            && \text{(Compound interest, continuous)} \\
  E_{1c} &: P(t) = P_0 \cdot \exp(k t)
            && \text{(Population growth)} \\
  E_{1d} &: C(t) = C_0 \cdot \exp(-k_{\mathrm{el}}\, t)
            && \text{(Pharmacokinetics, 1-compartment)}
\end{align*}

\paragraph{Explicit bijection $\varphi_{ab}$ for $E_{1a} \siso E_{1b}$.}

Label the nodes of $\MinDAG(E_{1a})$:
$v_1 = \mmul$ (root),\ $v_2 = \mexp$,\ $v_3 = \mmul$ (inner),\ $v_4 = N_0$
(terminal),\ $v_5 = -\lambda$ (terminal),\ $v_6 = t$ (terminal).

Label the nodes of $\MinDAG(E_{1b})$:
$u_1 = \mmul$ (root),\ $u_2 = \mexp$,\ $u_3 = \mmul$ (inner),\ $u_4 = V_0$,\
$u_5 = r$,\ $u_6 = t$.

Define $\varphi_{ab}(v_i) = u_i$ for $i = 1,\ldots,6$.  Verification:
\begin{itemize}
  \item \textit{Operator preservation:} Both $v_i$ and $u_i$ carry the same
        operator label for $i = 1,2,3$ and are terminals for $i = 4,5,6$.
  \item \textit{Topology:} $v_4, v_1 \to v_1$ (root), $v_2 \to v_1$,
        $v_3 \to v_2$, $v_5, v_6 \to v_3$ matches $u_4, u_1 \to u_1$, etc.
  \item \textit{Terminals map to terminals:} $\varphi_{ab}(\{v_4,v_5,v_6\})
        = \{u_4,u_5,u_6\}$.
\end{itemize}
Hence $E_{1a} \siso E_{1b}$.  By transitivity (Proposition~\ref{prop:equiv}),
all four members of Family~1 are mutually isomorphic.

\begin{remark}[Structural constant vs.\ domain parameter]
The inner multiplication $\mmul(b, c)$ in $T_1$ absorbs one sign flip:
radioactive decay uses $b = -\lambda < 0$ while growth uses $b = k > 0$.  The
\emph{structure} is identical; only the sign of a terminal value differs.
SuperBEST isomorphism intentionally ignores terminal values.
\end{remark}

% ── 2.2 Sum-of-exponentials family ───────────────────────────────────────────
\subsection{Family 2: Sum of Exponential-Decay Terms (NPV and N-Compartment PK)}
\label{sec:fam2}

\paragraph{Canonical tree.}
Let $T_1$ be as in~\eqref{eq:fam1-tree}.  For a sum of $N$ such terms, the
canonical tree is
\[
  T_2^{(N)} \;=\; \madd\bigl(T_1^{(1)},\; \madd(T_1^{(2)},\;\ldots\;
                  \madd(T_1^{(N-1)},\; T_1^{(N)})\ldots)\bigr),
\]
where $T_1^{(i)}$ is a fresh copy of $T_1$ with independent terminals
$(a_i, b_i, c_i)$.  Cost: $N \cdot 5 + (N-1) \cdot 3 = 8N - 3$
(using $c_{\madd} = 3$ in the positive domain, per T38 and the O($N$) theorem).

\paragraph{Member equations.}
\begin{align*}
  E_{2a} &: \mathrm{NPV}
    = \sum_{i=1}^{N} \frac{C_i}{(1+r)^i}
    = \sum_{i=1}^{N} C_i \cdot \exp\!\bigl(-i\,\ln(1+r)\bigr)
    && \text{(Finance)} \\
  E_{2b} &: \mathrm{AUC}
    = \sum_{i=1}^{N} A_i \cdot \exp(-\lambda_i t)
    && \text{(N-compartment PK)}
\end{align*}

Both are sums of $N$ terms, each term of the form $a_i \cdot \exp(b_i \cdot c_i)$
with $(a_i, b_i, c_i) = (C_i,\, -i,\, \ln(1+r))$ for NPV and
$(A_i,\, -\lambda_i,\, t)$ for multi-compartment PK.  The per-term canonical
tree is exactly $T_1$.

\paragraph{Explicit bijection for $E_{2a} \siso E_{2b}$ (one-term case).}

Within each summand $i$, the bijection is the same $\varphi_{ab}$ from
Family~1 with the substitution $N_0 \mapsto C_i$, $-\lambda \mapsto -i$,
$t \mapsto \ln(1+r)$ on the left and $V_0 \mapsto A_i$, $r \mapsto -\lambda_i$,
$t \mapsto t$ on the right.  Extending term-by-term and mapping each
$\madd$ node identically gives the full bijection for any $N \geq 1$.

\begin{corollary}
$E_{2a} \siso E_{2b}$ for all $N \geq 1$, with $\Cost(E_{2a}) =
\Cost(E_{2b}) = 8N - 3$.
\end{corollary}

% ── 2.3 Softmax / Logit ──────────────────────────────────────────────────────
\subsection{Family 3: Softmax and Logit Choice Probability}
\label{sec:fam3}

\paragraph{Canonical tree.}
\[
  T_3 \;=\; \mdiv\!\Bigl(
    \mexp(a),\;
    \madd\bigl(\mexp(b_1),\;\ldots,\;\madd(\mexp(b_{N-1}),\mexp(b_N))\bigr)
  \Bigr)
\]
with terminal slots $a, b_1, \ldots, b_N$.
Cost of numerator: $c_{\mexp} = 1$.
Cost of denominator ($N$-term sum of $\mexp$ nodes):
$N \cdot 1 + (N-1) \cdot 3 = 4N - 3$.
Total: $1 + 1 + (4N - 3) = 4N - 1$.
(The extra 1 comes from the $\mdiv$ node at the root.)

\paragraph{Member equations.}
\begin{align*}
  E_{3a} &: \sigma_i
    = \frac{\exp(x_i)}{\sum_{j=1}^{N} \exp(x_j)}
    && \text{(Softmax, machine learning)} \\
  E_{3b} &: P_i
    = \frac{\exp(v_i)}{\sum_{j=1}^{N} \exp(v_j)}
    && \text{(Logit / multinomial choice, economics)}
\end{align*}

\paragraph{Bijection.}
The two expressions are \emph{literally the same formula} with terminal
relabelling $x_j \mapsto v_j$ for $j = 1, \ldots, N$.  The identity map on
operator nodes, composed with the terminal relabelling, is the required
bijection.  Hence $E_{3a} \siso E_{3b}$ trivially, and their shared cost is
$4N - 1$.

\begin{remark}[Disciplinary distance]
Softmax originates in numerical optimisation for neural networks; logit choice
models originate in mathematical economics (McFadden's discrete choice theory).
They are separated by half a century of independent disciplinary development yet
their minimal operator trees are identical.
\end{remark}

% ── 2.4 Van't Hoff / Clausius-Clapeyron ──────────────────────────────────────
\subsection{Family 4: Van't Hoff Equation and Clausius-Clapeyron Relation}
\label{sec:fam4}

\paragraph{Canonical tree.}
\[
  T_4 \;=\; \madd\!\bigl(\mmul(a,\;\mrecip(T)),\;b\bigr)
\]
with terminal slots $a$ (slope), $T$ (temperature variable), $b$ (intercept).
Cost: $c_{\madd} = 3$, $c_{\mmul} = 2$, $c_{\mrecip} = 2$; total $= 7$.

\paragraph{Member equations.}
\begin{align*}
  E_{4a} &: \ln K(T)
    = -\frac{\Delta H^\circ}{R}\cdot\frac{1}{T} + \frac{\Delta S^\circ}{R}
    && \text{(Van't Hoff, thermodynamic equilibrium)} \\[4pt]
  E_{4b} &: \ln\frac{P_2}{P_1}
    = -\frac{\Delta H_{\mathrm{vap}}}{R}\!\left(\frac{1}{T_2}-\frac{1}{T_1}\right)
    && \text{(Clausius-Clapeyron, two-temperature form)}
\end{align*}

For $E_{4a}$: $a = -\Delta H^\circ/R$ (a single constant),
$T$ is the temperature, $b = \Delta S^\circ/R$.
For $E_{4b}$ (fixing $T_1$ as a compile-time constant):
$a = -\Delta H_{\mathrm{vap}}/R$, $T = T_2$,
$b = \Delta H_{\mathrm{vap}}/(RT_1)$.
Both reduce to $T_4$.

\paragraph{Explicit bijection $\varphi_{4}$ for $E_{4a} \siso E_{4b}$.}

Nodes of $\MinDAG(E_{4a})$:
$v_1 = \madd$ (root), $v_2 = \mmul$, $v_3 = \mrecip$,
$v_4 = a_{4a}$ (terminal $= -\Delta H^\circ/R$),
$v_5 = T$ (terminal), $v_6 = b_{4a}$ (terminal $= \Delta S^\circ/R$).

Nodes of $\MinDAG(E_{4b})$:
$u_1 = \madd$, $u_2 = \mmul$, $u_3 = \mrecip$,
$u_4 = a_{4b}$ (terminal $= -\Delta H_{\mathrm{vap}}/R$),
$u_5 = T_2$ (terminal), $u_6 = b_{4b}$ (terminal).

The map $\varphi_4(v_i) = u_i$ for $i=1,\ldots,6$ satisfies all three
conditions of Definition~\ref{def:iso}.  Hence $E_{4a} \siso E_{4b}$, with
$\Cost(E_{4a}) = \Cost(E_{4b}) = 7$.

% ── 2.5 Saturation / charging family ─────────────────────────────────────────
\subsection{Family 5: Saturation and Charging Curves \texorpdfstring{$(1 - e^{-t/\tau})$}{(1-exp(-t/tau))}}
\label{sec:fam5}

\paragraph{Canonical tree.}
\[
  T_5 \;=\; \msub\!\bigl(1,\; \mexp(\mmul(\mneg(1),\; \mdiv(t,\;\tau)))\bigr)
\]
with terminal slots $t$ and $\tau$.  The constants $1$ (unity) and $-1$ (or
the leading $\mneg$) are structural.  Cost: $c_{\msub}=2$, $c_{\mexp}=1$,
$c_{\mmul}=2$, $c_{\mneg}=2$, $c_{\mdiv}=1$; total $= 8$.  (When the full
expression is scaled by an amplitude $A$ there is an additional $\mmul$ root
node, giving cost $10$; the core saturation sub-tree has cost $8$.)

\paragraph{Member equations.}
\begin{align*}
  E_{5a} &: V(t) = V_0\!\bigl(1 - \exp(-t/RC)\bigr)
            && \text{(RC circuit charging)} \\
  E_{5b} &: T(t) = T_{\mathrm{env}}
              + (T_0 - T_{\mathrm{env}})\!\bigl(1 - \exp(-t/\tau)\bigr)
            && \text{(Newton's law of cooling)} \\
  E_{5c} &: y(t) = E_0\!\bigl(1 - \exp(-t/\tau)\bigr)
            && \text{(Thermal subsidence, geology)} \\
  E_{5d} &: C(t) = \frac{FD k_a}{V(k_a - k_e)}
              \bigl(\exp(-k_e t) - \exp(-k_a t)\bigr)
            && \text{(PK absorption, simplified)}
\end{align*}

Note: $E_{5b}$ has an additive offset $T_{\mathrm{env}}$ that adds a
$\madd$ root (cost $+3$), but the saturation core remains $T_5$.
$E_{5d}$ is a difference of two $T_1$-type terms with a shared constant
prefactor; its core saturation is expressed as $e^{-k_e t} - e^{-k_a t}$,
a rearrangement of $T_5$ via the identity
$1 - e^{-k_a t} - (1 - e^{-k_e t}) = e^{-k_e t} - e^{-k_a t}$.
The shared core structure is confirmed by equal cost of the saturation sub-tree.

\paragraph{Bijection $\varphi_{5ac}$ for $E_{5a} \siso E_{5c}$ (core trees).}
Label the six internal and terminal nodes of each core tree identically as in
Family 1 (substituting $RC \mapsto \tau$ and $V_0 \mapsto E_0$).  The node-wise
identity map on operators satisfies all three conditions.  Hence $E_{5a} \siso E_{5c}$.

% ── 2.6 Simple ratio family ───────────────────────────────────────────────────
\subsection{Family 6: Simple Ratio Laws (1n)}
\label{sec:fam6}

\paragraph{Canonical tree.}
\[
  T_6 \;=\; \mdiv(a,\; b)
\]
with two terminal slots.  Cost: $c_{\mdiv} = 1$.  This is the minimum
non-trivial cost achievable under $\mathcal{F}_{16}$.

\paragraph{Member equations.}
\begin{align*}
  E_{6a} &: \Phi = Q/\varepsilon_0       && \text{(Gauss flux law)} \\
  E_{6b} &: U = kq_1q_2/r               && \text{(Coulomb PE, fixed $kq_1q_2$)} \\
  E_{6c} &: \chi = C/T                  && \text{(Curie's law, magnetism)} \\
  E_{6d} &: v = d/t                     && \text{(Kinematics)} \\
  E_{6e} &: PV = C/r                    && \text{(Perpetuity, finance)}
\end{align*}

In each case the expression reduces to a single division of two quantities, one
of which may itself be a folded constant (e.g.\ $kq_1q_2$ in $E_{6b}$).
After constant folding, every member has the canonical form $\mdiv(a, b)$.

\paragraph{Bijection.}
The bijection is the identity on the single internal node $\mdiv$ and the
unique terminal-to-terminal map $a \mapsto a$, $b \mapsto b$ (renaming).
All five equations are mutually isomorphic, each with cost $1$.

\begin{remark}[1n equations as the degenerate class]
Family~6 sits at the absolute cost floor: $\Cost = 1$, the minimum for any
non-trivial expression.  The diversity of scientific domains (electrostatics,
magnetism, kinematics, finance) that produce 1n equations demonstrates that
the simplest ratio structure appears universally.
\end{remark}

% ── 2.7 Flux / transport laws ────────────────────────────────────────────────
\subsection{Family 7: Linear Transport Laws (Ohm, Darcy, Fick)}
\label{sec:fam7}

\paragraph{Canonical tree.}
\[
  T_7 \;=\; \mmul(G,\; \Delta V)
\]
equivalently written as $\mdiv(\Delta V,\; R)$ via the relation $G = 1/R$.
The division form costs $1$; the multiplication form costs $2$.  We adopt the
division form as canonical.

\paragraph{Member equations.}
\begin{align*}
  E_{7a} &: I = \Delta V / R
            && \text{(Ohm's law, electrical)} \\
  E_{7b} &: Q = -kA\,\Delta T / L
            && \text{(Darcy's law / Fourier heat, with sign absorbed)}\\
  E_{7c} &: J = -D\,\Delta c / \Delta x
            && \text{(Fick's first law of diffusion, with sign absorbed)}
\end{align*}

After absorbing the sign into the terminal constant (constant folding),
each reduces to $\mdiv(\Delta,\; R_{\mathrm{eff}})$ where $\Delta$ is the
driving potential difference and $R_{\mathrm{eff}}$ is the effective resistance.
Cost: $1$ in each case.

\paragraph{Bijection.}
Identical to Family~6: map the single $\mdiv$ node to itself and rename
terminals.  Hence $E_{7a} \siso E_{7b} \siso E_{7c}$, all with cost $1$.

\begin{remark}[Transport analogy as isomorphism]
The long-observed ``transport analogy'' in engineering (electrical, thermal,
diffusive, hydraulic circuits all obey the same $I = \Delta V / R$ template)
is exactly the statement that these equations lie in the same SuperBEST
isomorphism class.  The isomorphism \emph{is} the analogy, made precise.
\end{remark}

% ── 2.8 Entropy / information family ─────────────────────────────────────────
\subsection{Family 8: Entropy and Information (p~ln~p terms)}
\label{sec:fam8}

\paragraph{Canonical tree for a single term.}
\[
  T_8 \;=\; \mmul(p,\; \mln(p))
\]
with one terminal slot $p$.  Cost: $c_{\mmul} = 2$, $c_{\mln} = 1$; total $= 3$.
The full $N$-term sum costs $3N + 3(N-1) = 6N - 3$.

\paragraph{Member equations (per-term structure).}
\begin{align*}
  E_{8a} &: H = -\sum_{i=1}^{N} p_i \ln p_i
            && \text{(Shannon entropy, information theory)} \\
  E_{8b} &: S = -k_B \sum_{i=1}^{N} p_i \ln p_i
            && \text{(Boltzmann entropy, statistical physics)} \\
  E_{8c} &: D_{\mathrm{KL}} = \sum_{i=1}^{N} p_i \ln(p_i/q_i)
            && \text{(KL divergence, statistics)}
\end{align*}

For $E_{8a}$ and $E_{8b}$: each term is $p_i \cdot \ln(p_i) = T_8$.
The constant $-1$ (or $-k_B$) is a prefactor folded into the outer
$\mmul$ at the sum level.  For $E_{8c}$: each term is
$p_i \cdot \ln(p_i / q_i) = p_i \cdot (\ln p_i - \ln q_i)$; the per-term
tree is $\mmul(p_i, \msub(\mln(p_i), \mln(q_i)))$, with cost
$2 + 2 + 1 + 1 = 6$.  Hence $E_{8a} \siso E_{8b}$ (same per-term tree $T_8$),
while $E_{8c}$ lies in a distinct isomorphism class.

\paragraph{Bijection for $E_{8a} \siso E_{8b}$ (per-term).}
Map the single $\mmul$ node to itself and the single $\mln$ node to itself.
Map terminal $p_i$ to $p_i$ (relabelled; in $E_{8b}$ the same probability
appears with the Boltzmann constant absorbed into the outer sum).  The bijection
is the identity on operators plus terminal relabelling.

\begin{remark}[KL divergence as a different class]
The two-terminal-per-term structure of KL divergence ($p_i$ and $q_i$
independently appear) places it outside the isomorphism class of Shannon and
Boltzmann entropy, even though all three are ``entropy-like.''  SuperBEST
isomorphism distinguishes them at the level of the per-term DAG.
\end{remark}

% ════════════════════════════════════════════════════════════════════════════
\section{The SuperBEST Isomorphism Theorem}
\label{sec:theorem}
% ════════════════════════════════════════════════════════════════════════════

\begin{theorem}[SuperBEST Isomorphism Theorem]
\label{thm:main}
The eight families defined in Section~\ref{sec:families} are genuine
equivalence classes under $\siso$.  Specifically:
\begin{enumerate}
  \item[(F1)] All four exponential scale-and-decay equations (radioactive decay,
        compound interest, population growth, 1-compartment PK) are mutually
        SuperBEST-isomorphic, with $\Cost = 5$.
  \item[(F2)] The NPV sum and the $N$-compartment PK AUC sum are
        SuperBEST-isomorphic for every $N \geq 1$, with $\Cost = 8N - 3$.
  \item[(F3)] Softmax and logit choice probability are SuperBEST-isomorphic
        for every $N \geq 1$, with $\Cost = 4N - 1$.
  \item[(F4)] The Van't Hoff equation and the Clausius-Clapeyron relation
        (two-temperature form with one temperature fixed) are SuperBEST-isomorphic,
        with $\Cost = 7$.
  \item[(F5)] The RC charging curve, Newton cooling (core), and thermal
        subsidence are SuperBEST-isomorphic via the $1 - e^{-t/\tau}$ core,
        with core $\Cost = 8$.
  \item[(F6)] Gauss flux, Coulomb PE (fixed prefactor), Curie's law,
        kinematics $v = d/t$, and the perpetuity formula are mutually
        SuperBEST-isomorphic, with $\Cost = 1$.
  \item[(F7)] Ohm's law, Darcy's law (Fourier heat), and Fick's law of
        diffusion are mutually SuperBEST-isomorphic, with $\Cost = 1$.
  \item[(F8)] Shannon entropy and Boltzmann entropy are SuperBEST-isomorphic
        (per-term and full sum), with per-term $\Cost = 3$ and full-sum
        $\Cost = 6N - 3$.
\end{enumerate}
\end{theorem}

\begin{proof}
For each family we exhibit an explicit bijection satisfying
Definition~\ref{def:iso} in the corresponding subsection of
Section~\ref{sec:families}.  Mutual isomorphism of all members within a family
follows by transitivity (Proposition~\ref{prop:equiv}).  That the families are
genuine equivalence classes follows from the distinctness of their canonical
trees: the canonical trees of distinct families differ in at least one of
operator set, operator count, or topology, so no bijection satisfying conditions
(1)--(3) can exist between members of different families.
\end{proof}

\begin{table}[h]
\centering
\caption{Summary of the eight SuperBEST isomorphism families.}
\label{tab:families}
\smallskip
\begin{tabular}{@{}clL{5.5cm}c@{}}
\toprule
\textbf{Family} & \textbf{Canonical form} & \textbf{Members (domains)} & \textbf{Cost} \\
\midrule
F1 & $\mmul(a,\mexp(\mmul(b,c)))$ &
  Radioactive decay, compound interest, population growth, 1-cpt PK &
  5 \\[4pt]
F2 & $\sum_i \mmul(a_i,\mexp(\mmul(b_i,c_i)))$ &
  NPV (finance), $N$-cpt PK AUC &
  $8N\!-\!3$ \\[4pt]
F3 & $\mdiv(\mexp(a_i),\sum_j\mexp(b_j))$ &
  Softmax (ML), logit choice (economics) &
  $4N\!-\!1$ \\[4pt]
F4 & $\madd(\mmul(a,\mrecip(T)),b)$ &
  Van't Hoff, Clausius-Clapeyron &
  7 \\[4pt]
F5 & $\msub(1,\mexp(\mmul(\mneg(1),\mdiv(t,\tau))))$ &
  RC charging, Newton cooling, thermal subsidence &
  8 \\[4pt]
F6 & $\mdiv(a,b)$ &
  Gauss, Coulomb PE, Curie, kinematics, perpetuity &
  1 \\[4pt]
F7 & $\mdiv(\Delta V, R)$ &
  Ohm's law, Darcy/Fourier heat, Fick's diffusion &
  1 \\[4pt]
F8 & $\mmul(p,\mln(p))$ (per term) &
  Shannon entropy, Boltzmann entropy &
  $6N\!-\!3$ \\
\bottomrule
\end{tabular}
\end{table}

% ── 3.1 Uniqueness of canonical trees ────────────────────────────────────────
\subsection{Uniqueness of Canonical Representatives}

\begin{proposition}[Distinctness of family canonical trees]
\label{prop:distinct}
No two families in Table~\ref{tab:families} share the same canonical tree.
\end{proposition}

\begin{proof}
We distinguish the families pairwise by structural invariants:
\begin{itemize}
  \item F1 vs.\ F2: F2 has $N \geq 2$ summand copies of the F1 tree joined by
        $\madd$; F1 has $N = 1$ with no $\madd$ node.
  \item F1 vs.\ F3: F3 uses $\mdiv$ at the root and a sum in the denominator;
        F1 uses $\mmul$ at the root with no $\mdiv$.
  \item F4: the only family whose root is $\madd$ with one child being
        $\mmul(\cdot, \mrecip(\cdot))$ and whose cost is $7$ (odd).
  \item F5: the only family whose root is $\msub$ with right child $\mexp$.
  \item F6 and F7: both have canonical tree $\mdiv(a,b)$ and cost $1$.
        However, F6 and F7 \emph{do} share the same canonical tree; they are
        placed in separate families only for domain-annotation purposes.
        Formally $[E_{6a}] = [E_{7a}]$ under $\siso$.
  \item F8: the only family whose root is $\mmul$ with right child $\mln$.
\end{itemize}
With the note that F6 and F7 share a canonical tree (cost-1 division), all
other families are distinguished by at least one structural invariant.
\end{proof}

\begin{remark}[F6 and F7 as one class]
Formally, all ten equations in F6 and F7 belong to the \emph{same}
SuperBEST isomorphism class, the cost-1 division class.  We list them as
two families for expository clarity, grouping by the nature of the driving
quantity (fixed constant numerator in F6 vs.\ potential difference in F7).
The mathematics does not distinguish them.
\end{remark}

% ── 3.2 Corollary ────────────────────────────────────────────────────────────
\subsection{The Revealed Structure Corollary}

\begin{corollary}[Revealed Structure]
\label{cor:revealed}
For any two equations $E_1$, $E_2$ that are SuperBEST-isomorphic, there exists
a renaming of variables and constants in $E_1$ that transforms it into $E_2$
(after constant folding).  Equivalently, $E_1$ and $E_2$ are the same
computation, written in different notation.
\end{corollary}

\begin{proof}
By Definition~\ref{def:iso}, the bijection $\varphi$ establishes a one-to-one
correspondence between operators and topological roles.  Replacing each terminal
of $\MinDAG(E_1)$ with the label of the corresponding terminal in $\MinDAG(E_2)$
under $\varphi$ yields $\MinDAG(E_2)$ exactly, operator-for-operator.  This
is precisely the renaming claim.
\end{proof}

\begin{corollary}[Cross-domain transfer]
\label{cor:transfer}
Any analytic technique, approximation, or algorithm that applies to one member
of a SuperBEST isomorphism family applies, after terminal relabelling, to every
other member.
\end{corollary}

\begin{proof}
Any algorithm that operates on the operator structure of $\MinDAG(E_1)$
(without inspecting terminal labels) produces the same result when applied to
$\MinDAG(E_2)$ via the isomorphism $\varphi$.  Renaming terminals at the output
translates the result back to the domain of $E_2$.
\end{proof}

% ════════════════════════════════════════════════════════════════════════════
\section{Discussion: What Revealed Structure Means}
\label{sec:discussion}
% ════════════════════════════════════════════════════════════════════════════

% ── 4.1 Notation hides structure ─────────────────────────────────────────────
\subsection{How Domain Notation Conceals Isomorphism}

The equations in Family~1 look different in their native notation.  A physicist
writes $N(t) = N_0 e^{-\lambda t}$ in the context of decay chains and half-lives.
A financial analyst writes $V(t) = V_0 e^{rt}$ in the context of continuous
compounding and yield curves.  A pharmacokineticist writes
$C(t) = C_0 e^{-k_{\mathrm{el}} t}$ in the context of drug clearance.
The three communities rarely interact, and their textbooks do not cross-reference.

Yet at the level of the minimal operator tree, these three equations are
identical.  The SuperBEST DAG strips away the domain-specific labels and reveals
the shared skeleton: one outer multiplication, one exponential, one inner
multiplication, and three free terminals.  The isomorphism is not metaphorical;
it is a formal bijection between graphs.

% ── 4.2 Implications for the monogate system ────────────────────────────────
\subsection{Implications for monogate}

The Isomorphism Theorem has direct implications for the monogate routing system:

\begin{enumerate}
  \item \textbf{Cost prediction by class membership.}
        Identifying which isomorphism class an input expression belongs to
        immediately yields its cost (Table~\ref{tab:families}).  This requires
        no tree enumeration, only pattern matching against the canonical
        representatives.

  \item \textbf{Cross-domain template reuse.}
        If monogate has already built the optimal EML tree for one member of
        an isomorphism class, it can produce the optimal tree for any other
        member by relabelling terminals.  No re-optimisation is needed.

  \item \textbf{Structural search index.}
        Equations can be indexed by their isomorphism class rather than by
        domain keyword.  A query for ``sum of exponentials'' retrieves both
        NPV and multi-compartment PK AUC, because they share a canonical tree,
        regardless of the vocabulary used to describe them.
\end{enumerate}

% ── 4.3 Relation to existing analogies ──────────────────────────────────────
\subsection{Relation to Known Cross-Domain Analogies}

The transport analogy (Family~7) is centuries old: James Clerk Maxwell himself
drew the parallel between electrical conduction and heat flow.  Our result shows
that this analogy is not merely heuristic but is \emph{exact at the operator
level}: Ohm, Fourier, and Fick produce identical $\MinDAG$s.

The softmax-logit identification (Family~3) appears in the machine learning
literature as ``softmax is multinomial logistic regression'' but is usually
stated as an algebraic identity, not a structural one.  The isomorphism
theorem makes the identity structural: the two communities independently
arrived at the same operator graph.

The Van't Hoff -- Clausius-Clapeyron identification (Family~4) is less
well-known.  Both equations are standard results in physical chemistry taught
in separate chapters (equilibrium constants vs.\ phase transitions) with no
cross-reference.  The common $\ln = \mathrm{slope} \times (1/T) + \mathrm{intercept}$
skeleton is visible only after rewriting both in minimal operator form.

% ── 4.4 Limits of isomorphism ────────────────────────────────────────────────
\subsection{Limits of SuperBEST Isomorphism}

SuperBEST isomorphism is a \emph{structural} equivalence, not a semantic one.
Two isomorphic equations may behave very differently as physical models: they
may have different domains of validity, different stability properties, and
different empirical content.  The isomorphism says only that they share the
same minimal computation graph.

In particular, the sign of the terminal assigned to $b$ in the inner
$\mmul(b, c)$ of Family~1 determines whether the equation describes growth
($b > 0$) or decay ($b < 0$).  The structure is the same; the behaviour
differs.  The monogate routing system treats this as a terminal-value difference,
not a structural one, which is correct: the operator tree is identical.

% ── 4.5 Open questions ───────────────────────────────────────────────────────
\subsection{Open Questions}

\begin{enumerate}
  \item \textbf{Completeness.}  Do the eight families above cover all
        isomorphism classes that appear in the standard scientific corpus of
        50--100 equations, or are there additional classes not yet identified?

  \item \textbf{Mixed classes.}  Class~D equations (containing both
        \texttt{exp} and \texttt{ln}) may define isomorphism classes not
        reducible to the families above.  Characterising these classes is open.

  \item \textbf{Parameterised families.}  Families F2, F3, and F8 are
        parameterised by $N$.  For fixed $N$, they define a unique canonical
        tree.  Whether there exist families parameterised by structural
        parameters other than $N$ (e.g.\ depth, branching factor) is open.

  \item \textbf{Algorithmic recognition.}  Given an expression $E$, efficiently
        determining which isomorphism class it belongs to (graph isomorphism
        restricted to operator DAGs) is a well-posed algorithmic problem.
        For the fixed-size canonical trees in Table~\ref{tab:families}, linear
        time matching suffices.  For variable-$N$ families, efficient pattern
        matching algorithms remain to be implemented in monogate.
\end{enumerate}

% ════════════════════════════════════════════════════════════════════════════
\section*{Acknowledgments}
% ════════════════════════════════════════════════════════════════════════════

This document is part of the monogate cost theory sprint (sessions COST-1
through COST-9, April 2026).  The isomorphism perspective grew from the
observation in COST-4 that equations from different structural classes appeared
to share operator skeletons.  The formal definition and the eight-family
classification were developed in this session.

\bigskip
\noindent\textit{Almaguer, A.R.\ (2026).
The SuperBEST Isomorphism Theorem.
monogate research. \url{https://monogate.org}}

\end{document}
